\section{Présentation du projet}
\subsection{Contexte métier}
L'entreprise utilise un catalogue de prestations de travaux, de matériaux, de main-d'oeuvre et d'options techniques. Les demandes reçues peuvent être détaillées, partielles ou formulées oralement. Le projet vise à accélérer la création d'un brouillon de devis sans remplacer la validation du professionnel.

\subsection{Objectif général}
L'objectif est de concevoir une application capable de transformer une description de chantier en lignes de devis exploitables. La solution doit :
\begin{itemize}
\item reconnaître une ou plusieurs prestations dans une demande ;
\item extraire une quantité lorsqu'elle est explicitement donnée ou directement calculable ;
\item sélectionner uniquement les options et choix présents dans le catalogue ;
\item calculer les coûts et les prix de vente à partir des composants disponibles ;
\item présenter un aperçu clair, produire un PDF et préparer l'envoi au client ;
\item proposer des entrées texte, photo et vocales.
\end{itemize}

\subsection{Périmètre et limites}
Le système assiste la préparation du devis. Les montants, les prestations proposées et les quantités restent vérifiables et modifiables par le commercial. Une quantité absente de la demande n'est pas assimilée à zéro ; elle est conservée à \texttt{null} jusqu'à validation ou saisie complémentaire.
